\documentclass[10pt]{article}
\usepackage[T1]{fontenc}
\usepackage[utf8]{inputenc}
\usepackage[french]{babel}
\usepackage{amsmath,amssymb,mathrsfs,mathpazo}
\usepackage{geometry}
\geometry{a4paper,margin=2cm}
\usepackage{graphicx,array,tabularx,tabularray,colortbl,multirow,multicol,booktabs,enumitem,pifont,ulem,fancybox,diagbox,cancel}
\usepackage{pstricks,pstricks-add,pst-plot,pst-tree,pst-node,pst-text,pst-func,pst-math,pst-bspline,pst-3dplot,pst-eucl,pst-solides3d}
\let\pscalculate\undefined
\usepackage{pst-bezier}
\usepackage{tikz,pgfplots,tkz-tab}
\usetikzlibrary{arrows,arrows.meta,patterns,calc,positioning,trees,decorations.markings,shapes.geometric,shapes.misc,shapes.arrows,intersections,decorations.pathmorphing,decorations.pathreplacing,backgrounds,shadows,babel}
\pgfplotsset{compat=1.18}
\usepgfplotslibrary{fillbetween}
\usepackage[np]{numprint}
\usepackage{listings,makecell,pas-tableur,twoopt}
\let\barre\undefined
\usepackage[pstricks]{tablvar}
\definecolor{darkgreen}{RGB}{0,100,0}
\providecommand{\red}{\color{red}}
\providecommand{\blue}{\color{blue}}
\providecommand{\green}{\color{green}}
\providecommand{\black}{\color{black}}
\providecommand{\white}{\color{white}}
\providecommand{\gray}{\color{gray}}
\providecommand{\cyan}{\color{cyan}}
\providecommand{\magenta}{\color{magenta}}
\providecommand{\yellow}{\color{yellow}}
\providecommand{\violet}{\color{violet}}
\providecommand{\orange}{\color{orange}}
\providecommand{\euro}{\texteuro}
\providecommand{\EUR}{\texteuro}
\providecommand{\ieme}{\textsuperscript{e}}
\providecommand{\degres}{\ensuremath{^\circ}}
\providecommand{\degree}{\ensuremath{^\circ}}
\providecommand{\vv}[1]{\overrightarrow{#1}}
\providecommand{\Coord}[1]{#1}
\providecommand{\decosix}{}
\providecommand{\starredbullet}{\ensuremath{\bullet}}
\pagestyle{empty}
\setlength{\parindent}{0pt}
\usepackage{adjustbox}\DeclareUnicodeCharacter{2236}{:}\DeclareUnicodeCharacter{2190}{\ensuremath{\leftarrow}}
\providecommand{\R}{\mathbb{R}}
\providecommand{\N}{\mathbb{N}}
\providecommand{\Z}{\mathbb{Z}}
\providecommand{\Q}{\mathbb{Q}}
\providecommand{\D}{\mathbb{D}}
\providecommand{\C}{\mathbb{C}}
\providecommand{\e}{{\mathrm{e}}}
\providecommand{\vect}[1]{\overrightarrow{#1}}
\providecommand{\vectt}[1]{\overrightarrow{\mathrm{#1}}}
\providecommand{\vvt}[1]{\overrightarrow{\mathrm{#1}}}
\providecommand{\np}[2][]{{#2~#1}}
\providecommand{\vv}[1]{{\overrightarrow{#1}}}
\renewcommand{\d}{{\mathrm{d}}}
\providecommand{\degres}{{^\circ}}
\providecommand{\degre}{{^\circ}}
\providecommand{\blue}{\color{blue}}
\providecommand{\red}{\color{red}}
\providecommand{\enskip}{\quad}
\providecommand{\llbracket}{\lbrack\!\lbrack}
\providecommand{\rrbracket}{\rbrack\!\rbrack}
\providecommand{\euro}{{\text{€}}}
\providecommand{\no}{{\text{n°}}}
\providecommand{\interval}[2]{{[#1;#2]}}
\providecommand{\intervFO}[2]{{[#1;#2[}}
\providecommand{\dotfill}{\ldots}

\providecommand{\ds}{\displaystyle}
\providecommand{\textding}[1]{#1}
\providecommand{\Oij}{\(\mathrm{O};\vec{\imath},\vec{\jmath}\)}
\providecommand{\Oijk}{\(\mathrm{O};\vec{\imath},\vec{\jmath},\vec{k}\)}
\providecommand{\Ouv}{\(\mathrm{O};\vec{u},\vec{v}\)}
\providecommand{\intervFF}[2]{\ensuremath{[#1;#2]}}
\renewcommand{\euro}{\ifmmode\text{\texteuro}\else\texteuro\fi}
\DeclareUnicodeCharacter{00B0}{\ifmmode^{\circ}\else\textdegree\fi}
\DeclareUnicodeCharacter{20AC}{\ifmmode\text{\texteuro}\else\texteuro\fi}
% Énoncés : Éducation nationale. Transcription et corrigés : APMEP, Denis Vergès et contributeurs. Découpage et classement : CoopMaths, Sébastien Lozano, Éric Elter et contributeurs. Adaptation de la présentation : Maths Explorer.
% CC BY-SA : https://coopmaths.fr/www/about/
% https://coopmaths.fr/alea/static/e3c/2020/tex/e3c_2020_00_sujet13_2_cor.tex
\begin{document}
\noindent\textit{APMEP / CoopMaths. CC BY-SA. Présentation adaptée par Maths Explorer.}\par\medskip
\subsection*{1. Montrer que la suite \((u_n)\) est une suite géométrique dont on donnera la raison.}
Ajouter 3,7 \% c'est multiplier par \(1 + \dfrac{3,7}{100} = 1 + 0,037 = 1,037\).\\
 On a donc pour tout naturel \(n\), \(u_{n+1} = 1,037u_n\) : la suite \((u_n)\) est une suite géométrique de premier terme \(u_0 = 187\) et de raison 1,037.
	
\subsection*{2. Pour tout \(n \in \mathbb{N}\), exprimer \(u_n\) en fonction de \(n\).}
On sait qu'alors avec \(q\) comme raison \(u_n = u_0 \times q^n\), quel que soit le naturel \(n\). Ici \(u_n = 187 \times 1,037^n\).
	
\subsection*{3. Étudier le sens de variation de la suite \((u_n)\).}
	
	
	
	La suite \((u_n)\) est géométrique de premier terme \(u_0=187\) et de raison \(q=1,037\).
	On a \(u_0 > 0\) et \(q>1\), la suite \((u_n)\) est donc strictement croissante.
	
	\subsection*{4. Selon cette estimation, calculer la production mondiale de plastique en 2019. Arrondir au million de tonnes.}
	
	2019 correspond à \(n = 19\), \\
	donc \(u_{19} = 187 \times 1,037^{19} \approx 372,9\), donc environ 373 millions au million près.
	
	\subsection*{5. Des études montrent que 20 \% de la quantité totale de plastique se retrouve dans les océans, et que 70 \% de ces déchets finissent par couler. Montrer que la quantité totale, arrondie au million de tonnes, de déchets flottants sur l'océan dus à la production de plastique de 2000 à 2019 compris est de 324 millions de tonnes.}
	
	Sur les 20 \% de plastiques allant à la mer, 30 \% flottent en 2000, soit 6 \%. La production totale de plastiques de 2000 à 2019 est :
	\[ S_{2019} = u_0 + 1,037u_0 + \ldots + 1,037^{19}u_0 \]
	\[ 1,037S_{2019} = 1,037u_0 + 1,037^2u_0 + \ldots + 1,037^{20}u_0 \]
	Par différence :
	\[ 0,037S_{2019} = 1,037^{20}u_0 - u_0 \]
	\[ S_{2019} = u_0 \dfrac{1,037^{20} - 1}{0,037} \approx 5398,32 \]
	Restent en surface :
	\[ 5398,32 \times 0,20 \times 0,30 \approx 323,89, \]
	soit 324 millions au million près.

\end{document}
